\scanpage{407}\mspage{146}
famille à deux indices de sous-catégories
pleines
\[
\underline{\mathrm{Ch}}_{n,m}^{\deleted{\unsure{W}}}(X)\qquad n,m\in\mathbb Z,\quad\deleted{\unsure{n\leq m}}
\]
en définissant cette catégorie comme \uline{vide} si $n>m$,
et pour $n\leq m$ comme formée des objets
\diagram{\(\left\{\begin{gathered}
\deleted{\unsure{\biggl\{}}(a_i;u_i)_{i\in\mathbb Z}\in\mathrm{Ch}^W_{n,m}(X)\Longleftrightarrow
\begin{aligned}&u_i\text{ une identité pour}\\& i\notin[n,m-1]\end{aligned}\\[5pt]
\begin{tikzcd}[ampersand replacement=\&,column sep=2pt,font=\scriptsize,execute at end picture={
\draw (an.south) .. controls +(0,-.16) and +(-.18,.15) .. ($(an.south)!.5!(am.south)+(0,-.27)$) .. controls +(.18,.15) and +(0,-.16) .. (am.south);
\node[below,font=\scriptsize,align=center] at ($(an.south)!.5!(am.south)+(0,-.28)$) {``\unsure{portion|partie} utile''\\du chemin};}]
\overset{\mathrm{id}}{\cdots}\&a_{n-1}\arrow[r,no head,"\mathrm{id}"]\&|[alias=an]|a_n\arrow[r,no head,"u_n"]\&a_{n+1}\&\cdots\arrow[r,no head]\&a_{m-1}\arrow[r,no head,"u_{m-1}"]\&|[alias=am]|a_m\arrow[r,no head,"\mathrm{id}"]\&a_{m+1}\&\overset{\mathrm{id}}{\cdots}
\end{tikzcd}
\end{gathered}\right.\)}

Donc pour $n=m$, on trouve la sous-catégorie pleine
isomorphe à $X$, formée des chemins infinis
``totalement dégénérés'' i.e. tels que \uline{toutes} les
$u_i$ soient des identités. (Cette sous-cat. de
$\underline{\mathrm{Ch}}_{\infty}$ ne dépend donc pas du choix de $n=m$.)
\diagram{\begin{tikzcd}[ampersand replacement=\&,row sep=tiny,column sep=small]
{}\arrow[r,"\mathrm{id}"]\&a_{-1}\arrow[d,phantom,sloped,"=" description]\&a_0\arrow[l,"\mathrm{id}"']\arrow[r,"\mathrm{id}"]\arrow[d,phantom,sloped,"=" description]\&a_1\arrow[d,phantom,sloped,"=" description]\&{}\arrow[l,"\mathrm{id}"']\\
\&a\&\unsure{a|c}\&a\&
\end{tikzcd}}
Tous les sommets d'un \deleted{\unsure{même}} \inserted{tel} chemin sont égaux
à un même objet \deleted{\unsure{$a$}} de $X$, et les flèches
de transition entre eux sont identiques.

La structure de $\underline{\mathrm{Ch}}^W_{n,m}(X)\simeq\underline{\mathrm{Ch}}^W_{\tau_{n,m}}(X)$
ne dépend que de la longueur $m-n$ de $\tau_{n,m}$,
et de la \uline{\unsure{parité}} de $n$. Si $n$ est pair,
la direction initiale de $\tau_{n,m}$ est positive,
et
\[
\tau_{n,m}=\underbrace{\Delta_+\Delta_-\ \ldots}_{\substack{m-n\ \text{facteurs}\\\text{en alternance}}}
=\begin{cases}
(\Delta_+\Delta_-)^{\frac{m-n}{2}}&\text{si }m\text{ pair}\\
(\Delta_+\Delta_-)^{\frac{m-n}{2}-1}\Delta_+&\text{si }\inserted{\unsure{m\text{ impair}|\text{en général}}}
\end{cases}
\]
tandis que pour $n$ impair c'est juste l'inverse
% archive sheet number: 406
